
Un rayo luminoso, \textit{AB}, propagándose en el aire, incide en una de las caras de una lámina de vidrio de caras paralelas. Después de sufrir dos refracciones, el rayo de luz emerge de la lámina para el aire (rayo \textit{CD}), como se muestra en la figura de este problema.
\begin{enumerate}
    \item Aplique la ley de Snell para determinar la relación entre los ángulos $\theta_1$ y $\theta_4$ mostrados en la figura.
    \item Teniendo en cuenta la respuesta (\textit{a}), ¿qué relación existe entre las direcciones de los rayos \textit{AB} y \textit{CD}?
\end{enumerate}